% !PRINTED 364
\weberpage{364}
%<<CONT
zustellen. Wir beschränken uns hier auf das, was für die Erreichung unseres Hauptzieles, nämlich der Bestimmung der \gesperrt{Galois}'schen Gruppe des Problems, erforderlich ist.

Wir gehen von einem beliebig herausgegriffenen \gesperrt{Steiner}'schen Complexe aus:
\weq{1}{x_1 y_1,\ x_2 y_2,\ x_3 y_3,\ x_4 y_4,\ x_5 y_5,\ x_6 y_6.}
Die Form der Gleichung \S.\,96, (5) zeigt dann, dass in dem durch das Paar $x_1 x_2$ bestimmten Complexe das Paar $y_1 y_2$, und in dem durch $x_1 y_2$ bestimmten Complexe das Paar $x_2 y_1$ vorkommen muss.

Wir betrachten also neben (1) die zwei weiteren Complexe:
\weq{2}{x_1 x_2,\quad y_1 y_2,\quad \ldots}
\weq{3}{x_1 y_2,\quad x_2 y_1,\quad \ldots}
Jeder der beiden Complexe (2), (3) enthält ausser den schon bekannten noch acht weitere Doppeltangenten, und diese müssen alle von den $x_i, y_i$ verschieden sein, weil, wenn z. B. $x_3$ in (2) vorkäme, $x_1 x_2 x_3$ syzygetisch wäre, was dem Satze 3., (\S.\,96) widerspricht. Ebenso können die beiden Complexe (2) und (3) ausser $x_1, x_2, y_1, y_2$ keine gemeinsame Doppeltangente enthalten. Denn wenn etwa $z$ in beiden vorkäme, so wären $x_1 x_2 z$ nach (2) syzygetisch, nach (3) azygetisch, was ein Widerspruch ist. Daraus folgt:

\gesperrt{4. In den drei Complexen (1), (2), (3) zusammengenommen kommen alle 28 Doppeltangenten vor.}

Hieraus folgt, dass jede Doppeltangente, die mit irgend einem Paare $x_1 y_1$ ein syzygetisches Tripel bildet, in dem Complexe $x_1 y_1$ vorkommen muss, dass also jedes syzygetische Tripel durch eine bestimmte weitere Doppeltangente zu einem syzygetischen Quadrupel ergänzt wird, und dass folglich ein Kegelschnitt, der durch die Berührungspunkte von drei Doppeltangenten hindurchgeht, die Curve $f$ in den Berührungspunkten einer vierten Doppeltangente schneidet.

Zwei Paare eines Complexes bilden immer ein syzygetisches Quadrupel, und wie man ein solches Quadrupel auch in zwei Paare theilen mag, beide Paare gehören immer demselben Complexe an.

Da man aus 28 Dingen 14.27 Paare bilden kann, da jedes Paar von Doppeltangenten einen Complex bestimmt und in jedem
%CONT>>
